\chapter*{Введение}
\addcontentsline{toc}{chapter}{Введение}

Малые и средние предприятия часто принимают управленческие решения на основе разрозненных Excel-выгрузок, банковских выписок и данных учетных систем. В результате собственники могут видеть общий оборот бизнеса, но не иметь прозрачного ответа на более важные вопросы: какие категории расходов съедают прибыль, почему возникают кассовые разрывы, насколько бизнес зависит от отдельных покупателей и поставщиков, какие операции требуют ручной проверки.

Тема данной курсовой работы связана с применением методов машинного обучения и больших языковых моделей для финансового и операционного анализа малого производственного бизнеса. В качестве практического примера рассматривается компания, производящая кондитерские изделия, печенье и выпечку. Компания продает продукцию через сети дискаунтеров, магазины низких цен и торговые рынки Москвы, Московской области и Тульской области.

Актуальность работы определяется тем, что даже небольшой бизнес накапливает значительный объем транзакционных данных. Эти данные уже содержат информацию о продажах, закупках, возвратах, аренде, налогах, кредитах и платежах поставщикам, однако без очистки и классификации они остаются бухгалтерским массивом строк, а не инструментом управления.

Цель работы -- разработать и проверить подход к применению методов машинного обучения и LLM для анализа транзакционных и операционных данных малого производственного бизнеса, чтобы повысить прозрачность денежных потоков, выявить ключевые финансовые риски и подготовить основу для управленческих решений.

Для достижения цели поставлены следующие задачи:
\begin{enumerate}
  \item описать специфику применения ML и LLM в малом и среднем бизнесе;
  \item провести обзор источников данных компании;
  \item построить воспроизводимый пайплайн очистки и подготовки данных;
  \item разработать схему классификации банковских операций;
  \item сравнить rule-based и LLM-подходы к разметке транзакций;
  \item сформировать управленческий cash flow;
  \item провести финансовый аудит денежных потоков;
  \item обучить и оценить прогнозные модели;
  \item подготовить рекомендации по внедрению результатов в управленческую практику.
\end{enumerate}

Объект исследования -- малое производственное предприятие в сфере кондитерских изделий. Предмет исследования -- методы автоматической обработки, классификации и прогнозирования транзакционных и операционных данных предприятия.

Практическая значимость работы состоит в том, что предложенный подход переводит разовую работу с выгрузками в воспроизводимый контур: данные очищаются, операции классифицируются, денежный поток пересчитывается, модели прогнозирования обновляются, а спорные операции выделяются для ручного контроля.
